% !TEX program = xelatex

%-------------------------------------%
% Birmingham City University
% Joint Program with University of Information Technology, VNU-HCM
% Module: Individual Honours Project
%
% -----
%
% A Game-Based Approach to Motor Imagery Adaptation and BCI Calibration
% Author: Cao Đăng Khoa (dang.cao@mail.bcu.ac.uk)
% Student ID: 25195675
%
% Supervisor: Dr. Vi Chí Thành (vcthanh@hcmiu.edu.vn)
%------------------------------------%



%------ DOCUMENT CONFIGURATION ------%
\documentclass[a4paper, 11pt, openany]{report}
\usepackage{layout}
\usepackage{fonts}
\usepackage{headerfooter}
\usepackage{frontmatter}
\usepackage{contents}
\usepackage{atbegshi}

\makeatletter
\let\ps@plain\ps@fancy
\makeatother
%------------------------------------%



%----------- FRONT MATTER -----------%
\begin{document}

% Title Page
\reportTitlePage
    {Project Interim Report}
    {A Game-Based Approach to Motor Imagery Adaptation and BCI Calibration}
    {Cao Đăng Khoa}{25195675}
    {Bachelor of Science with Honours Computer Science}
    {Dr. Vi Chí Thành}{December 1, 2025}
    \thispagestyle{titlepage}

% Page numbering
\pagenumbering{roman}
\setcounter{page}{0}

% Table of Contents
\reportTableOfContents

% List of Figures
\reportListOfFigures

% List of Tables
\reportListOfTables
%------------------------------------%


\cleardoublepage
\pagenumbering{arabic}

\setstretch{1.2} 
\setlength{\parskip}{8pt plus 2pt minus 1pt} 


%------------ CHAPTER 1 -------------%
%----- Introduction and Context -----%
\vspace*{-1.15cm}
\chapter{Introduction and Context}
    % 1.0 Introduction and Context
    This section introduces your project. It should be brief yet engaging, capturing the reader's 
    interest and highlighting the significance of your work. \par

    Begin broadly by introducing the subject and its main discipline, such as cybersecurity, 
    sustainable engineering, or healthcare analytics. Provide an overview of the background and 
    current situation, offering enough context so a reader unfamiliar with the field can understand 
    its significance. Why is this area of modern relevance? \par

    Note: Remember to use in-text citations for statements that support your core arguments. This 
    applies to all sections throughout your project report. 

    % 1.1 Problem Statement
    \section{Problem Statement}
    Clearly and specifically state the problem your project is designed to address. This should be 
    a concise statement that identifies the key issue. Think of it as the "what" and "why" of your 
    project, all in one or two sentences. For example: "Small to medium-sized enterprises (SMEs) 
    lack affordable and scalable cybersecurity solutions, leaving them highly susceptible to data 
    breaches." \par
%------------------------------------%



%------------ CHAPTER 2 -------------%
%--------- Literature Review --------%
\chapter{Review of Existing Knowledge}
    % 2.0 Review of Existing Knowledge
    Comprehensively outline and expand on the primary concepts and established models presented in
     studies relevant to your problem statement. What is the existing knowledge on this subject? What 
     are the main findings and conclusions derived from these studies? This should not merely be a 
     descriptive summary; rather, a critical analysis of the information to construct a coherent narrative 
     regarding the current state of knowledge. 
    \par 
    
    Note: The review of concepts could be broken down into further sub-sections depending on the themes 
    relevant to your topic that you identify during the review. 

    % 2.1 Critical Analysis and Gap Identification
    \section{Critical Analysis and Gap Identification}
    During this phase, you carry out a critical reflection of the sources you have examined, identifying 
    their limitations, contradictions, or disagreements within the literature. Establish which questions 
    remain unanswered. This critical scrutiny should help you pinpoint a specific "gap”; whether it is an 
    unresolved issue, an underexplored area, a new perspective that has not been considered, or a practical 
    solution that is missing from existing studies.

    % 2.2 Project Justification
    \section{Project Justification}
    Motor imagery-based BCIs are among the most commonly studied paradigms due to their non-invasive nature 
    and potential for rehabilitation applications.
%------------------------------------%



%------------- CHAPTER 3 -------------%
%-Project Aims, Objectives, and Scope-%
\chapter{Project Aims, Objectives, and Scope}
    % 3.0 Project Aims, Objectives, and Scope
    This section clearly defines what your project intends to achieve (aims), the specific steps you will 
    take to reach those goals (objectives), and the boundaries within which your project will operate (scope). 
    \par

    % 3.1 Project Aims
    \section{Project Aims}
    The primary aim of this project is to develop a game-based motor imagery BCI system that can adapt to
        individual users, thereby enhancing calibration efficiency and overall user experience.

    % 3.2 Project Objectives
    \section{Project Objectives}
    To achieve the stated aim, the project will pursue the following objectives:
\section{Project Scope}

\newpage
\chapter{ba asndjsadnjasdn}
\end{document}
